% GENERATED FILE. Edit canonical lessons, then run npm run generate:books.
% canonical-source-hash: fnv1a64:c48f1582e215241c
% canonical-lessons: LA-C334-sacerdos, LA-C334-victima, LA-C334-immolo, LA-C334-voveo, LA-C334-votum

\chapter{Sacrificium}
\label{ch:la-sacrificium}
\hlchaptermodality{334}

\begin{chapteropening}
\textbf{By the end of this chapter:} \emph{I can describe a sacrifice: the priest, the victim, to sacrifice, to vow and a vow.}

Complete the last lesson of chapter 334: I can describe a sacrifice: the priest, the victim, to sacrifice, to vow and a vow.
\end{chapteropening}

\section[sacerdōs, sacerdōtis]{sacerdōs, sacerdōtis — a priest, a priestess}

\label{lesson:LA-C334-sacerdos}



\begin{quote}
Before the new one: say the Latin for a feast, then the Latin for to feast.
\end{quote}

\subsection*{You'll want to know: sacerdōs, sacerdōtis}
\textbf{sacerdōs, sacerdōtis} — \textquotedblleft{}a priest, a priestess\textquotedblright{}.

\begin{quote}
Sacerdōs ante āram stābat.
\end{quote}

\begin{quote}
\emph{The priest was standing before the altar.}
\end{quote}

Masculine or feminine, 3rd declension, from \emph{sacer}, 'holy'.

\begin{cousinweb}
\textbf{sacerdōs} is built on \emph{sacer}, \textquotedblleft{}holy\textquotedblright{} (EtymDB), with a second part whose origin is not settled: one who performs holy things.
\end{cousinweb}

\begin{grammarlens}[title={What you've built}]
One more word for telling, explaining and describing at B1.
\end{grammarlens}

\subsection*{Your turn}
\begin{itemize}
\raggedright
  \item \emph{Say it:} \emph{sacerdōs, sacerdōtis}
  \item \emph{Say it:} \emph{sacerdōs, sacerdōtis}, once more
  \item \emph{Say it:} a sentence of your own with \emph{sacerdōs, sacerdōtis}
  \item \emph{Recall:} say the Latin for a feast, then the Latin for to feast, then say \emph{sacerdōs, sacerdōtis} again
\end{itemize}

\begin{tcolorbox}[breakable,colback=teal!4,colframe=teal!35!black,title={Before you move on}]
What does sacerdōs, sacerdōtis mean? (\textquotedblleft{}A priest, a priestess\textquotedblright{}.) Say it once more.
\end{tcolorbox}
\section[victima, victimae]{victima, victimae — a sacrificial animal, a victim}

\label{lesson:LA-C334-victima}



\begin{quote}
Before the new one: say the Latin for to feast, then the Latin for a priest.
\end{quote}

\subsection*{You'll want to know: victima, victimae}
\textbf{victima, victimae} — \textquotedblleft{}a sacrificial animal, a victim\textquotedblright{}.

\begin{quote}
Victima alba ad āram ducta est.
\end{quote}

\begin{quote}
\emph{A white victim was led to the altar.}
\end{quote}

Feminine. A larger animal for sacrifice, such as an ox. English \emph{victim}.

\begin{cousinweb}
\textbf{victima} is a Latin word whose origin is not settled; Lewis \& Short suggest, with doubt, a root seen in \emph{vigeō}, \textquotedblleft{}to be strong\textquotedblright{}. English \emph{victim} comes from it.
\end{cousinweb}

\begin{grammarlens}[title={What you've built}]
One more word for telling, explaining and describing at B1.
\end{grammarlens}

\subsection*{Your turn}
\begin{itemize}
\raggedright
  \item \emph{Say it:} \emph{victima, victimae}
  \item \emph{Say it:} \emph{victima, victimae}, once more
  \item \emph{Say it:} a sentence of your own with \emph{victima, victimae}
  \item \emph{Recall:} say the Latin for to feast, then the Latin for a priest, then say \emph{victima, victimae} again
\end{itemize}

\begin{tcolorbox}[breakable,colback=teal!4,colframe=teal!35!black,title={Before you move on}]
What does victima, victimae mean? (\textquotedblleft{}A sacrificial animal, a victim\textquotedblright{}.) Say it once more.
\end{tcolorbox}
\section[immolō, immolāre]{immolō, immolāre — to sacrifice}

\label{lesson:LA-C334-immolo}



\begin{quote}
Before the new one: say the Latin for a priest, then the Latin for a sacrificial animal.
\end{quote}

\subsection*{You'll want to know: immolō, immolāre}
\textbf{immolō, immolāre} — \textquotedblleft{}to sacrifice\textquotedblright{}.

\begin{quote}
Sacerdōs taurum deae immolāvit.
\end{quote}

\begin{quote}
\emph{The priest sacrificed a bull to the goddess.}
\end{quote}

\emph{immolō, immolāre, immolāvī, immolātum}: first, to sprinkle the victim with salted grain.

\begin{cousinweb}
\textbf{immolō} is built on \emph{mola}, \textquotedblleft{}meal, salted grain\textquotedblright{}, with \emph{in-} (Lewis \& Short): to sprinkle the victim with sacred meal before the sacrifice.
\end{cousinweb}

\begin{grammarlens}[title={What you've built}]
One more word for telling, explaining and describing at B1.
\end{grammarlens}

\subsection*{Your turn}
\begin{itemize}
\raggedright
  \item \emph{Say it:} \emph{immolō, immolāre}
  \item \emph{Say it:} \emph{immolō, immolāre}, once more
  \item \emph{Say it:} a sentence of your own with \emph{immolō, immolāre}
  \item \emph{Recall:} say the Latin for a priest, then the Latin for a sacrificial animal, then say \emph{immolō, immolāre} again
\end{itemize}

\begin{tcolorbox}[breakable,colback=teal!4,colframe=teal!35!black,title={Before you move on}]
What does immolō, immolāre mean? (\textquotedblleft{}To sacrifice\textquotedblright{}.) Say it once more.
\end{tcolorbox}
\section[voveō, vovēre]{voveō, vovēre — to vow, to promise solemnly}

\label{lesson:LA-C334-voveo}



\begin{quote}
Before the new one: say the Latin for a sacrificial animal, then the Latin for to sacrifice.
\end{quote}

\subsection*{You'll want to know: voveō, vovēre}
\textbf{voveō, vovēre} — \textquotedblleft{}to vow, to promise solemnly\textquotedblright{}.

\begin{quote}
Dux deae templum novum vōvit.
\end{quote}

\begin{quote}
\emph{The general vowed a new temple to the goddess.}
\end{quote}

\emph{voveō, vovēre}, perfect \emph{vōvī}: short \emph{o} in the present, long in the perfect and in its noun.

\begin{cousinweb}
\textbf{voveō} is inherited; EtymDB traces it to a Proto-Indo-European root meaning \textquotedblleft{}to vow\textquotedblright{}, though Lewis \& Short call the origin doubtful. English \emph{vow} and \emph{vote} come from its noun \emph{vōtum}.
\end{cousinweb}

\begin{grammarlens}[title={What you've built}]
One more word for telling, explaining and describing at B1.
\end{grammarlens}

\subsection*{Your turn}
\begin{itemize}
\raggedright
  \item \emph{Say it:} \emph{voveō, vovēre}
  \item \emph{Say it:} \emph{voveō, vovēre}, once more
  \item \emph{Say it:} a sentence of your own with \emph{voveō, vovēre}
  \item \emph{Recall:} say the Latin for a sacrificial animal, then the Latin for to sacrifice, then say \emph{voveō, vovēre} again
\end{itemize}

\begin{tcolorbox}[breakable,colback=teal!4,colframe=teal!35!black,title={Before you move on}]
What does voveō, vovēre mean? (\textquotedblleft{}To vow, to promise solemnly\textquotedblright{}.) Say it once more.
\end{tcolorbox}
\section[vōtum, vōtī]{vōtum, vōtī — a vow; a wish}

\label{lesson:LA-C334-votum}



\begin{quote}
Before the new one: say the Latin for to sacrifice, then the Latin for to vow.
\end{quote}

\subsection*{You'll want to know: vōtum, vōtī}
\textbf{vōtum, vōtī} — \textquotedblleft{}a vow; a wish\textquotedblright{}.

\begin{quote}
Post naufragium nautae vōtum solvērunt.
\end{quote}

\begin{quote}
\emph{After the shipwreck the sailors fulfilled their vow.}
\end{quote}

Neuter, from \emph{voveō}. \emph{Vōtum solvere}: to keep a vow. English \emph{vote} and \emph{vow}.

\begin{cousinweb}
\textbf{vōtum} is built on \emph{voveō}, \textquotedblleft{}to vow\textquotedblright{} (EtymDB): a thing vowed. English \emph{vote} comes from it, and \emph{vow} through French.
\end{cousinweb}

\begin{grammarlens}[title={What you've built}]
One more word for telling, explaining and describing at B1.
\end{grammarlens}

\subsection*{Your turn}
\begin{itemize}
\raggedright
  \item \emph{Say it:} \emph{vōtum, vōtī}
  \item \emph{Say it:} \emph{vōtum, vōtī}, once more
  \item \emph{Say it:} a sentence of your own with \emph{vōtum, vōtī}
  \item \emph{Recall:} say the Latin for to sacrifice, then the Latin for to vow, then say \emph{vōtum, vōtī} again
\end{itemize}

\begin{tcolorbox}[breakable,colback=teal!4,colframe=teal!35!black,title={Before you move on}]
What does vōtum, vōtī mean? (\textquotedblleft{}A vow; a wish\textquotedblright{}.) Say it once more.
\end{tcolorbox}
